\Needspace{8\baselineskip}
\SourceLeaf{320}
\Needspace{8\baselineskip}
\Article{277}
Die Regeln, mit denen alle Formen in I und II des vorigen Artikels bestimmt werden können, folgen unmittelbar aus ihrer Beschreibung; einige Beispiele werden daher genügen. Für \(D=1\) ergeben sich durch die verschiedenen Vorzeichen diese sechs Formen I.
\[
\left(\begin{smallmatrix}0,&1,&0\\0,&\pm1,&0\end{smallmatrix}\right),\quad
\left(\begin{smallmatrix}0,&1,&\pm1\\0,&\pm1,&0\end{smallmatrix}\right)
\]
Bei den Formen II können \(a\) und \(A''\) keine anderen Werte als \(+1\) und \(-1\) haben. Für jede der vier daraus entstehenden Kombinationen müssen \(b''\), \(B\) und \(B'\) \(=0\) gesetzt werden; daraus entstehen vier Formen
\[
\left(\begin{smallmatrix}1,&-1,&1\\0,&0,&0\end{smallmatrix}\right),\quad
\left(\begin{smallmatrix}-1,&1,&1\\0,&0,&0\end{smallmatrix}\right),\quad
\left(\begin{smallmatrix}1,&1,&-1\\0,&0,&0\end{smallmatrix}\right),\quad
\left(\begin{smallmatrix}-1,&-1,&-1\\0,&0,&0\end{smallmatrix}\right)
\]
Ebenso erhält man für \(D=-1\) sechs Formen I und vier Formen II,
\[
\begin{gathered}
\left(\begin{smallmatrix}0,&-1,&0\\0,&\pm1,&0\end{smallmatrix}\right),\quad
\left(\begin{smallmatrix}0,&-1,&\pm1\\0,&\pm1,&0\end{smallmatrix}\right);\quad
\left(\begin{smallmatrix}1,&-1,&-1\\0,&0,&0\end{smallmatrix}\right),\\
\left(\begin{smallmatrix}-1,&1,&-1\\0,&0,&0\end{smallmatrix}\right),\quad
\left(\begin{smallmatrix}-1,&-1,&1\\0,&0,&0\end{smallmatrix}\right),\quad
\left(\begin{smallmatrix}1,&1,&1\\0,&0,&0\end{smallmatrix}\right)
\end{gathered}
\]
Für \(D=2\) entstehen sechs Formen I
\[
\left(\begin{smallmatrix}0,&2,&0\\0,&\pm1,&0\end{smallmatrix}\right),\quad
\left(\begin{smallmatrix}0,&2,&\pm1\\0,&\pm1,&0\end{smallmatrix}\right)
\]
und acht Formen II
\[
\begin{gathered}
\left(\begin{smallmatrix}1,&-1,&2\\0,&0,&0\end{smallmatrix}\right),\quad
\left(\begin{smallmatrix}-1,&1,&2\\0,&0,&0\end{smallmatrix}\right),\quad
\left(\begin{smallmatrix}1,&1,&-2\\0,&0,&0\end{smallmatrix}\right),\quad
\left(\begin{smallmatrix}-1,&-1,&-2\\0,&0,&0\end{smallmatrix}\right)\\
\left(\begin{smallmatrix}1,&-2,&1\\0,&0,&0\end{smallmatrix}\right),\quad
\left(\begin{smallmatrix}-1,&2,&1\\0,&0,&0\end{smallmatrix}\right),\quad
\left(\begin{smallmatrix}1,&2,&-1\\0,&0,&0\end{smallmatrix}\right),\quad
\left(\begin{smallmatrix}-1,&-2,&-1\\0,&0,&0\end{smallmatrix}\right)
\end{gathered}
\]
Die Anzahl der Klassen, die sich in diesen drei Fällen aus diesen Formen ergeben, ist jedoch viel kleiner als die Anzahl der Formen. Es lässt sich nämlich leicht bestätigen:

I. Die Form \(\left(\begin{smallmatrix}0,&1,&0\\0,&1,&0\end{smallmatrix}\right)\) geht in
\[
\left(\begin{smallmatrix}0,&1,&0\\0,&-1,&0\end{smallmatrix}\right),\quad
\left(\begin{smallmatrix}0,&1,&1\\0,&\pm1,&0\end{smallmatrix}\right),\quad
\left(\begin{smallmatrix}0,&1,&-1\\0,&\pm1,&0\end{smallmatrix}\right),\quad
\left(\begin{smallmatrix}1,&1,&-1\\0,&0,&0\end{smallmatrix}\right)
\]
jeweils durch die Substitutionen
\[
\begin{array}{rrr|rrr|rrr|rrr}
1,&0,&0&0,&0,&1&0,&0,&1&1,&0,&-1\\
0,&1,&0&0,&1,&-1&0,&1,&1&1,&1,&-1\\
0,&0,&-1&\pm1,&1,&0&\pm1,&-1,&-1&0,&-1,&1
\end{array}
\]
über, die Form \(\left(\begin{smallmatrix}1,&1,&-1\\0,&0,&0\end{smallmatrix}\right)\) dagegen in \(\left(\begin{smallmatrix}1,&-1,&1\\0,&0,&0\end{smallmatrix}\right)\), \(\left(\begin{smallmatrix}-1,&1,&1\\0,&0,&0\end{smallmatrix}\right)\) allein durch Vertauschung der Unbestimmten. Jene zehn ternären Formen mit der Determinante \(1\) reduzieren sich also auf diese beiden
\SourceLeaf{321}
\(\left(\begin{smallmatrix}0,&1,&0\\0,&1,&0\end{smallmatrix}\right)\), \(\left(\begin{smallmatrix}-1,&-1,&-1\\0,&0,&0\end{smallmatrix}\right)\); für die erste kann man, wenn man möchte, auch \(\left(\begin{smallmatrix}1,&0,&0\\1,&0,&0\end{smallmatrix}\right)\) nehmen. Da die erstere Form indefinit und die letztere definit ist, ist klar, dass jede indefinite ternäre Form mit der Determinante \(1\) zur Form \(xx+2yz\) und jede definite zu \(-xx-yy-zz\) äquivalent ist.

II. Auf genau dieselbe Weise findet man, dass jede indefinite ternäre Form mit der Determinante \(-1\) zur Form \(-xx+2yz\) und jede definite zu \(xx+yy+zz\) äquivalent ist.

III. Für die Determinante \(2\) können von den acht Formen (II) sofort die zweite, sechste und siebte verworfen werden, da sie aus der ersten allein durch Vertauschung der Unbestimmten entstehen. Aus demselben Grund können auch die fünfte, die ebenso aus der dritten, und die achte, die aus der vierten entsteht, verworfen werden. Die drei übrigen bilden mit den sechs Formen I drei Klassen. Denn \(\left(\begin{smallmatrix}0,&2,&0\\0,&1,&0\end{smallmatrix}\right)\) geht in \(\left(\begin{smallmatrix}0,&2,&0\\0,&-1,&0\end{smallmatrix}\right)\) durch die Substitution \(\left\{\begin{smallmatrix}1,&0,&0\\0,&1,&0\\0,&0,&-1\end{smallmatrix}\right\}\) über, und die Form \(\left(\begin{smallmatrix}1,&1,&-2\\0,&0,&0\end{smallmatrix}\right)\) in
\[
\begin{gathered}
\left(\begin{smallmatrix}0,&2,&1\\0,&1,&0\end{smallmatrix}\right),\quad
\left(\begin{smallmatrix}0,&2,&1\\0,&-1,&0\end{smallmatrix}\right),\quad
\left(\begin{smallmatrix}0,&2,&-1\\0,&1,&0\end{smallmatrix}\right),\\
\left(\begin{smallmatrix}0,&2,&-1\\0,&-1,&0\end{smallmatrix}\right),\quad
\left(\begin{smallmatrix}1,&-1,&2\\0,&0,&0\end{smallmatrix}\right)
\end{gathered}
\]
jeweils durch die Substitutionen
\[
\begin{array}{rrr|rrr|rrr|rrr|rrr}
1,&0,&1&1,&0,&-1&1,&0,&0&1,&0,&0&1,&0,&0\\
1,&2,&0&1,&2,&0&1,&2,&-1&1,&2,&1&0,&1,&2\\
1,&1,&0&1,&1,&0&1,&1,&-1&1,&1,&1&0,&1,&1
\end{array}
\]
über. Jede ternäre Form mit der Determinante \(2\) lässt sich also auf eine dieser drei reduzieren
\[
\left(\begin{smallmatrix}0,&2,&0\\0,&-1,&0\end{smallmatrix}\right),\quad
\left(\begin{smallmatrix}1,&1,&-2\\0,&0,&0\end{smallmatrix}\right),\quad
\left(\begin{smallmatrix}-1,&-1,&-2\\0,&0,&0\end{smallmatrix}\right)
\]
Anstelle der ersten kann man, wenn man möchte, auch \(\left(\begin{smallmatrix}2,&0,&0\\1,&0,&0\end{smallmatrix}\right)\) nehmen. Offenbar muss jede definite ternäre Form zur dritten, \(-xx-yy-2zz\), äquivalent sein, da die beiden ersten indefinit sind. Jede indefinite ist zur ersten oder zweiten äquivalent: zur ersten, \(2xx+2yz\), wenn ihr erster, zweiter und dritter Koeffizient alle gerade sind (denn man sieht leicht, dass eine solche Form durch jede Substitution wieder in eine Form derselben Art übergeht und daher nicht zur zweiten Form äquivalent sein kann); zur zweiten, \(xx+yy-2zz\), dagegen, wenn ihr erster, zweiter und dritter Koeffizient nicht alle gerade, sondern einer, zwei oder alle ungerade sind (in eine solche Form kann nämlich aus demselben Grund die erste Form \(2xx+2yz\) durch keine Substitution verwandelt werden).
\SourceLeaf{322}
Was also in den Beispielen der Artikel 273, 274 geschah, dass die definite Form
$\left(\begin{smallmatrix}19,21,50\\15,28,1\end{smallmatrix}\right)$
mit der Determinante $-1$ auf $xx+yy+zz$ und die indefinite Form
$\left(\begin{smallmatrix}10,26,2\\7,0,4\end{smallmatrix}\right)$
mit der Determinante $2$ auf $2xx-2yz$ oder, was auf dasselbe hinausläuft, auf $2xx+2yz$ reduziert wurde, hätte man aus den vorigen Untersuchungen im Voraus erkennen können.

\Needspace{8\baselineskip}
\Article{278}
Durch eine ternäre Form mit den Unbestimmten $x,x',x''$ werden sowohl Zahlen \textit{dargestellt}, indem man $x,x',x''$ bestimmte Werte gibt, als auch binäre Formen durch Substitutionen der Art
\[
x=mt+nu,\qquad x'=m't+n'u,\qquad x''=m''t+n''u
\]
wobei $m,n,m'$ usw.\ bestimmte Zahlen und $t,u$ die Unbestimmten der dargestellten Form bezeichnen. Eine vollständige Theorie der ternären Formen würde daher die Lösung folgender Probleme erfordern: I.\ Alle Darstellungen einer gegebenen Zahl durch eine gegebene ternäre Form finden. II.\ Alle Darstellungen einer gegebenen binären Form durch eine gegebene ternäre Form finden. III.\ Entscheiden, ob zwei gegebene ternäre Formen mit derselben Determinante äquivalent sind, und wenn ja, alle Transformationen der einen in die andere finden. IV.\ Entscheiden, ob eine gegebene ternäre Form eine andere gegebene Form mit größerer Determinante enthält, und wenn ja, alle Transformationen jener in diese angeben. Diese Probleme, die viel schwieriger sind als die entsprechenden bei den binären Formen, werden wir an anderer Stelle ausführlicher behandeln. Hier beschränken wir unsere Untersuchung darauf zu zeigen, wie sich das erste Problem auf das zweite und das zweite auf das dritte zurückführen lässt. Für einige sehr einfache Fälle, die besonders die Theorie der binären Formen erhellen, werden wir die Lösung des dritten zeigen. Das vierte schließen wir hier ganz aus.

\Needspace{8\baselineskip}
\Article{279}
\textsc{Lemma.} \textit{Zu drei beliebigen gegebenen ganzen Zahlen $a,a',a''$ (die aber nicht alle $=0$ sind): sechs andere $B,B',B'',C,C',C''$ finden, sodass gilt}
\[
B'C''-B''C'=a,\qquad B''C-BC''=a',\qquad BC'-B'C=a''
\]

\textit{Lösung.} Es sei $\alpha$ der größte gemeinsame Teiler von $a,a',a''$. und man wähle ganze Zahlen $A,A',A''$ so, dass
\[
Aa+A'a'+A''a''=\alpha
\]
\SourceLeaf{323}
Ferner wähle man drei ganze Zahlen $\mathfrak{C},\allowbreak \mathfrak{C}',\allowbreak \mathfrak{C}''$ beliebig mit der einzigen Bedingung, dass die drei Zahlen
$\mathfrak{C}'A''-\mathfrak{C}''A',\allowbreak \ \mathfrak{C}''A-\mathfrak{C}A'',\allowbreak \
\mathfrak{C}A'-\mathfrak{C}'A$, die wir jeweils mit $b,b',b''$ und deren größten gemeinsamen Teiler wir mit $\beta$ bezeichnen werden, nicht alle $=0$ sind. Dann setze man
\[
a'b''-a''b'=\alpha\beta C,\qquad
a''b-ab''=\alpha\beta C',\qquad
ab'-a'b=\alpha\beta C''
\]
Offenbar sind $C,C',C''$ ganze Zahlen. Wählt man schließlich ganze Zahlen
$\mathfrak{B},\allowbreak \mathfrak{B}',\allowbreak \mathfrak{B}''$ so, dass
\[
\mathfrak{B}b+\mathfrak{B}'b'+\mathfrak{B}''b''=\beta
\]
setzt
\[
\mathfrak{B}a+\mathfrak{B}'a'+\mathfrak{B}''a''=h
\]
und bestimmt
\[
B=\alpha\mathfrak{B}-hA,\qquad
B'=\alpha\mathfrak{B}'-hA',\qquad
B''=\alpha\mathfrak{B}''-hA''
\]
so erfüllen diese Werte von $B,B',B'',C,C',C''$ die geforderten Gleichungen.

Denn man findet
\[
\begin{gathered}
aB+a'B'+a''B''=0\\
bA+b'A'+b''A''=0\quad\text{also}\quad bB+b'B'+b''B''=\alpha\beta
\end{gathered}
\]
Aus den Werten von $C',C''$ ergibt sich nun
\[
\begin{aligned}
\alpha\beta(B'C''-B''C')
&=ab'B'-a'bB'-a''bB''+ab''B''\\
&=a(bB+b'B'+b''B'')-b(aB+a'B'+a''B'')=\alpha\beta a
\end{aligned}
\]
also $B'C''-B''C'=a$; ebenso findet man $B''C-BC''=a'$, $BC'-B'C=a''$. \textit{Q.E.F.} --- Die Analyse, durch die diese Lösung gefunden wurde, sowie die Methode, aus einer Lösung alle zu finden, müssen wir hier weglassen.
